\chapter*{Abstract}
\addcontentsline{toc}{chapter}{Abstract}
MAYO is an Oil-and-Vinegar multivariate signature scheme submitted to post-quantum standardisation. In the official AVX2 C implementation, products of the public matrix $P_1$ by the vectors of $V$ are computed with \code{vpshufb} lookups on packed $\mathbb{F}_{16}$ vectors. When the vector length $m$ is not a multiple of 64, the official kernel covers the tail with an overlapping 32-byte load that duplicates work. We study a local change, called A2, which handles the 8-byte tail of each $P_1$ entry in a single 128-bit register and one xmm \code{vpshufb} on the low lane of the multiplication table, reducing the number of shuffles per scalar pair from 4 to 3 (MAYO-1) and from 6 to 5 (MAYO-5). We formalise the counting argument and its limits, establish the properties on which correctness rests, and measure the effect on a physical laptop with a paired protocol. Kernel gains are 12--29\,\%; the full-signature gain is 10.3\,\% for MAYO-1 but only 1.3\,\% for MAYO-5. Static analysis shows that in the signature context GCC~16.2.0 emits extra 32-byte stores in several loops, asymmetrically between the two versions, consistently with the measured loss; causality is not established. GCC's dumps show that the extra stores are per-iteration stores of the accumulator arrays that are not moved out of the loop; SRA is ruled out, and why the compiler makes this choice remains open.


\chapter*{Résumé}
\addcontentsline{toc}{chapter}{Résumé}
MAYO est un schéma de signature numérique « multivarié » de type Oil-and-Vinegar. Il est candidat à la standardisation de la cryptographie post-quantique. Dans le code C officiel (MAYO-C) pour AVX2, le programme multiplie la matrice publique $P_1$ par les vecteurs de $V$ avec l'instruction \code{vpshufb}. Les vecteurs de $\mathbb{F}_{16}$ y sont rangés à deux éléments par octet. Quand la longueur $m$ n'est pas un multiple de 64, le code officiel traite la fin du vecteur en relisant 32 octets déjà traités. Cela fait des calculs en double : 48 éléments sur 128 pour MAYO-1 ($m=80$) et 50 sur 192 pour MAYO-5 ($m=142$).

Ce mémoire étudie une petite modification, appelée A2. Elle traite la fin de 8 octets de chaque entrée de $P_1$ dans un seul registre de 128 bits : un seul \code{vpshufb} xmm remplace deux \code{vpshufb} ymm. Le nombre de \code{vpshufb} par paire de nombres passe de 4 à 3 (MAYO-1) et de 6 à 5 (MAYO-5). Nous démontrons cette limite de comptage, nous établissons les propriétés qui rendent A2 correct, puis nous mesurons l'effet sur un ordinateur portable (Intel Core i9-9880H, GCC~16.2.0). Les mesures comparent les deux versions dans les mêmes conditions, et l'unité statistique est le processus (une exécution du programme).

Voici les résultats établis. Sur les noyaux pris seuls, A2 réduit le temps de 19 à 22\,\% (\code{P1_times_Vt}) et de 12 à 29\,\% (\code{P1t_times_Vt}). Sur la signature complète, le gain est de 10,3\,\% pour MAYO-1, mais de seulement 1,3\,\% pour MAYO-5. Pour MAYO-5, il monte à 7,3\,\% quand toutes les sondes de mesure sont actives. En lisant le code produit par GCC~16.2.0 (sans l'exécuter), nous voyons que, sans toutes les sondes, le compilateur ajoute dans plusieurs boucles une écriture mémoire de 32 octets par accumulateur, et que ce n'est pas la même chose dans les deux versions. Les temps mesurés varient de la même façon, mais le lien de cause à effet n'est pas démontré, et un modèle à un seul paramètre exagère le gain d'environ 40\,\%. Ces écritures en plus sont les stockages des tableaux d'accumulateurs, refaits à chaque tour de boucle au lieu d'être faits une seule fois après la boucle. Nous écartons la passe SRA du compilateur comme cause. Pourquoi le compilateur fait ce choix reste une question ouverte.

\medskip
\noindent\textbf{Mots-clés :} MAYO, signatures multivariées, cryptographie post-quantique, AVX2, \code{vpshufb}, $\mathbb{F}_{16}$, optimisation, mesure de performance.
